% Source PDF page 119; manual page 2-90.
\subsection{Flight Path Limits}\label{sec:flight-path-limits}

In addition to guidance rules, constraints can be imposed on the flight path. These
limits are usually applied directly to control variables as upper and lower bounds for
the Newton--Raphson solution, although any guidance rule can also serve as a
flight-path limit.

\taossubhead{Limits on Specified Control Variables}

A limit applied to a directly specified control variable is ignored. For example, if
the guidance rule commands a $25^\circ$ angle of attack while a limit specifies
$15^\circ$, TAOS retains the commanded $25^\circ$ value.

\taossubhead{Limits on Free Control Variables}

A limit on a free control variable is enforced during the Newton--Raphson search.
Default bounds are $\pm30^\circ$ for angle of attack and sideslip and
$\pm180^\circ$ for the other attitude angles.

If the guidance solution fails to converge, reducing the free-control limits to a more
realistic range can help. Aerodynamic tables may not extrapolate smoothly to angles
as large as $\pm30^\circ$, and narrower bounds can keep the search in a region where
the table data are better behaved.

\taossubhead{Limits on Flight Conditions}

When a limit is applied to an indirect flight-condition rule, TAOS compares the
trajectory value with the limit at each integration step. Once the limit is exceeded,
the limiting rule replaces one of the existing guidance rules.

Each indirect rule has a prioritized list of control variables likely to satisfy it. TAOS
selects the highest-priority associated control and replaces the guidance rule
currently using that control with the limiting rule. For example, if a vehicle climbing
under angle-of-attack guidance reaches an altitude limit, the altitude rule replaces
the angle-of-attack rule and commands constant altitude at the limiting value.
